\section{Trabajo Futuro}
\label{sec:trabajo-futuro}

Se declara explícitamente lo que quedó fuera de alcance de este informe, en vez de omitirlo o
maquillarlo con contenido de relleno:

\begin{enumerate}[nosep]
    \item \textbf{Despliegue público en producción:} preparado completamente (Dockerfiles verificados,
    configuración de Railway, documentación de despliegue), pero la cuenta y el despliegue real requieren
    una acción del equipo humano que no se completó a la fecha de este informe.
    \item \textbf{Cobertura de pruebas automatizadas:} 36 de 50 requisitos Must tienen prueba automatizada
    dedicada (72\,\%, ver Sección~\ref{sec:requisitos}); los 14 restantes se enumeran en \S10.3 del SRS.
    Cobertura general de 71.46\,\% de líneas y 53.82\,\% de ramas (JaCoCo, 2026-09-13, 657 tests, 0
    fallos), la primera por encima del objetivo de 60\,\% de RNF-04, la segunda bajo el objetivo de
    70\,\%. La capa de controladores supera el umbral de 70\,\% exigido por la guía en líneas \emph{y} en
    ramas a la vez, con margen real: pasó de 8.75\,\%/0\,\% a \textbf{79.01\,\%/80.34\,\%}. Servicios pasó de
    59.64\,\%/40.87\,\% a
    \textbf{70.22\,\%/54.48\,\%}: supera el umbral en líneas pero no en ramas. La brecha de ramas que
    queda a nivel global se concentra en los componentes transversales de seguridad (autenticación, JWT,
    filtros) y en la propia capa de servicios, y queda declarada como trabajo futuro en vez de omitida.
    \item \textbf{Aplicación real del instrumento SUS:} el cuestionario está listo; aplicarlo a un grupo
    real de al menos 5--10 usuarios (estudiantes, docentes, coordinadores) es la tarea pendiente más
    directa y de menor esfuerzo entre las listadas aquí.
    \item \textbf{Corrección de las 14 vulnerabilidades de dependencias de build restantes en el
    frontend} (5 críticas, en \texttt{sharp}/\texttt{to-ico}/\texttt{jimp}, herramientas de generación
    de íconos que no se envían al navegador) --- la vulnerabilidad crítica que sí afectaba código
    servido a producción (\texttt{@angular/core}, detectada por ZAP) se corrigió el 2026-08-29.
    \item \textbf{Revisión sistemática de literatura formal:} la Sección~\ref{sec:trabajos-relacionados}
    es una búsqueda dirigida, no una SLR completa con protocolo pre-registrado y doble revisor
    independiente, siguiendo Kitchenham y Charters \cite{kitchenham2007}.
    \item \textbf{Persistencia de archivos subidos en producción:} el filesystem del contenedor backend
    es efímero en el despliegue de Railway planeado; los PDF de anteproyectos y actas se perderían en
    cada redeploy sin un volumen persistente o almacenamiento de objetos externo.
    \item \textbf{Firma real del docente-director sobre el SRS v1.0.0:} la página de aprobación está
    lista; la firma en sí es una acción pendiente del docente, no del equipo de desarrollo.
    \item \textbf{Ampliar las pruebas de integración reales:} el proyecto ya tiene dos clases de prueba
    de integración verdadera contra PostgreSQL (\texttt{@SpringBootTest} +
    \texttt{@AutoConfigureTestDatabase(replace = NONE)}, sin mocking del repositorio):
    \texttt{PreSustentacionesApplicationTests} y \texttt{TemaPropuestoRepositoryIntegrationTest} --- esta
    última fue precisamente la que detectó un bug real de inferencia de tipos SQL invisible a las
    pruebas unitarias mockeadas (Sección~\ref{sec:seguridad-corregida}), demostrando el valor del
    patrón. Sigue pendiente extenderlo a más repositorios/servicios con lógica de negocio compleja, en
    vez de depender solo de pruebas unitarias y \textit{slice tests} \texttt{@WebMvcTest} para ellos.
\end{enumerate}

Esta lista se prioriza deliberadamente: los ítems 1--3 son alcanzables en días con el equipo humano
actuando (no requieren investigación adicional, solo ejecución de lo ya preparado); los ítems 4--8
requieren más tiempo de desarrollo o coordinación externa (el docente-director, en el caso del ítem 7).
\newpage
